\scanpage{33}
\mspage{7)}

Nous allons maintenant examiner de plus près le cas où \deleted{$H^+$ \unsure{invarie}} $\mathrm{Herm}^{++}(V)^{H^+}$ n'est pas réduit à une demi-droite. On vérifie aisément que cela revient à dire que $H^+$ \deleted{\unsure{est tel}} fixe deux pts distincts $z,z'$ de $X$, donc est contenu dans le sous groupe de $G^+$ \deleted{défini} formé des pts qui fixent $z,z^{\unsure{'|*}}$, lequel est isom à $\mathbb C^*$, \deleted{\unsure{et il a}} \inserted{et admet un} unique sous-groupe cpt maximal \deleted{$\unsure{Y}$ \unsure{isom. à}} \inserted{$T$ isom à $\mathbb U$}. \deleted{Dans ce cas Les métriques invariées par}

Si $H^+\ne\{1\}$, \inserted{[alors $\{z,z'\}$ donc $T$ est uniquement déterminé]} les métriques invariées par $H^+$ sont aussi celles invariées par $T\supset H^+$, \margin{\uline{à vérifier}} et elles sont en corr. $1$-$1$ avec les \deleted{trajectoires} \inserted{orbites} de $T$ dans $X^*=X\smallsetminus\{z,z'\}$, [NB $T$ opère librement sur $X^*\deleted{\unsure{\setminus\{z}}$].

La relation est obtenue en associant à tte métrique invariée par $T$, \deleted{l'équateur} le grand cercle ayant comme pôles $z,z'$, lequel est une \deleted{trajectoire} \inserted{orbite} sous $T$. En particulier, tte orbite \inserted{de $X^*$} sous $H^+$ (étant contenue dans une unique orbite de $T$) définit une telle métrique.

Si $H=H^+$, il n'y a \deleted{\unsure{pas}} rien à ajouter. Dans le

% archive sheet number: 35
